\newpage
\section{绪论}
\subsection{机器学习：我们在学什么}

设有一组样本
\[
x_1,x_2,\dots,x_N,\qquad x_i\in\mathbb{R}^d.
\]
也就是说，每个样本 \(x_i\) 都是一个 \(d\) 维向量。这些样本可以看作是从某个真实但未知的数据分布中抽取出来的，这个分布记作\(P_{\mathrm{data}}(x)\)。

在机器学习中，我们通常希望构造一个带参数的模型分布\(P_\theta(x)\)，其中 \(\theta\) 表示模型参数。学习的目标就是调节参数 \(\theta\)，使得模型分布尽可能接近真实数据分布，即
\[
P_\theta(x)\approx P_{\mathrm{data}}(x).
\]

一个典型例子是高斯分布。若模型假设数据服从某个高斯分布，则其密度可写成
\[
p_\theta(x)\propto \exp\left(-\frac{(x-\mu)^2}{2\sigma^2}\right),
\]
参数为\(\theta=(\mu,\sigma)\)，而学习就是通过数据去确定合适的 \(\mu,\sigma\)，使这个高斯分布尽量接近真实数据的分布。

\subsection{如何衡量``接近''：K-L散度}

要让
\[
P_\theta(x)\approx P_{\mathrm{data}}(x),
\]
首先需要一个``距离''或``差异''的度量。这里引入一个非常重要的量：\textbf{K-L 散度}。

对于两个概率分布 \(P\) 和 \(Q\)，其 K-L 散度定义为
\[
D_{\mathrm{KL}}(P\|Q)\equiv\int p(x)\log\frac{p(x)}{q(x)}\,\mathrm{d}x.
\]
这里\(p(x),q(x)\)分别是分布 \(P,Q\) 的密度。它衡量的是：如果真实分布是 \(P\)，但我们用 \(Q\) 来近似，会有多大的信息损失。

把 K-L 散度展开：
\[
D_{\mathrm{KL}}(P\|Q)
=
\int p(x)\log p(x)\,\mathrm{d}x
-
\int p(x)\log q(x)\,\mathrm{d}x.
\]
注意到
\[
-\int p(x)\log p(x)\,\mathrm{d}x,
\]
这正是分布 \(P\) 的\textbf{熵}：
\[
H(P)=-\int p(x)\log p(x)\,\mathrm{d}x.
\]
通常定义\textbf{交叉熵}为
\[
H(P,Q)\equiv-\int p(x)\log q(x)\,\mathrm{d}x,
\]
于是有
\[
D_{\mathrm{KL}}(P\|Q)=H(P,Q)-H(P).
\]
这说明：
\begin{itemize}[leftmargin=2.5em]
    \item \(H(P)\) 是真实分布本身的不确定性；
    \item \(H(P,Q)\) 是用 \(Q\) 来描述 \(P\) 时的平均代价；
    \item 二者之差就是额外损失，也就是 K-L 散度。
\end{itemize}

\subsection{K-L散度的非负性}

K-L 散度最重要的性质之一是
\[
D_{\mathrm{KL}}(P\|Q)\ge 0.
\]
并且只有当\(P=Q\)才有\(D_{\mathrm{KL}}(P\|Q)=0\)。这说明 K-L 散度虽然不一定是严格意义上的``距离''，但至少满足非负且相同则为零这一重要性质，因此非常适合作为分布拟合的目标函数。

下面来推导这一性质。从定义出发，
\[
D_{\mathrm{KL}}(P\|Q)=\int p(x)\log\frac{p(x)}{q(x)}\,\mathrm{d}x.
\]
把它改写成
\[
D_{\mathrm{KL}}(P\|Q)=-\int p(x)\log\frac{q(x)}{p(x)}\,\mathrm{d}x.
\]

现在考虑函数
\[
f(t)=-\log t.
\]
这个函数是一个凸函数，因此可以使用 Jensen 不等式：
\[
\int p(x)f\left(\frac{q(x)}{p(x)}\right)dx
\ge
f\left(\int p(x)\frac{q(x)}{p(x)}dx\right).
\]
代入 \(f(t)=-\log t\)，得到
\[
-\int p(x)\log\frac{q(x)}{p(x)}\,\mathrm{d}x
\ge
-\log\left(\int p(x)\frac{q(x)}{p(x)}\,\mathrm{d}x\right).
\]
而右边括号内可化简为
\[
\int p(x)\frac{q(x)}{p(x)}\,\mathrm{d}x
=
\int q(x)\,\mathrm{d}x
=1.
\]
所以
\[
D_{\mathrm{KL}}(P\|Q)\ge -\log 1=0.
\]
因此得到
\[
D_{\mathrm{KL}}(P\|Q)\ge 0.
\]

\subsection{学习目标：MLE}

在概率建模中，我们希望用参数化模型分布\(P_\theta(x)\)去逼近真实的数据分布\(P_{\mathrm{data}}(x)\)，即选择参数 \(\theta\)，使得模型分布与真实分布之间的 K-L 散度最小：
\[
\min_\theta D_{\mathrm{KL}}(P_{\mathrm{data}}\|P_\theta).
\]
这里
\[
D_{\mathrm{KL}}(P_{\mathrm{data}}\|P_\theta)
=
\int p_{\mathrm{data}}(x)\log\frac{p_{\mathrm{data}}(x)}{p_\theta(x)}\,\mathrm{d}x.
\]

展开K-L散度
\[
D_{\mathrm{KL}}(P_{\mathrm{data}}\|P_\theta)
=
\int p_{\mathrm{data}}(x)\log p_{\mathrm{data}}(x)\,\mathrm{d}x
-
\int p_{\mathrm{data}}(x)\log p_\theta(x)\,\mathrm{d}x.
\]
注意到第一项
\[
\int p_{\mathrm{data}}(x)\log p_{\mathrm{data}}(x)\,\mathrm{d}x
\]
只和真实分布 \(P_{\mathrm{data}}\) 有关，与参数 \(\theta\) 无关。因此，在对 \(\theta\) 做优化时，这一项可以视为常数。于是最小化 K-L 散度就等价于最小化
\[
-\int p_{\mathrm{data}}(x)\log p_\theta(x)\,\mathrm{d}x,
\]
也等价于最大化
\[
\int p_{\mathrm{data}}(x)\log p_\theta(x)\,\mathrm{d}x.
\]
所以有
\[
\min_\theta D_{\mathrm{KL}}(P_{\mathrm{data}}\|P_\theta)
\quad\Longleftrightarrow\quad
\max_\theta \int p_{\mathrm{data}}(x)\log p_\theta(x)\,\mathrm{d}x.
\]

进一步，上式中的
\[
\int p_{\mathrm{data}}(x)\log p_\theta(x)\,\mathrm{d}x
\]
其实就是在真实数据分布下，\(\log p_\theta(x)\) 的期望：
\[
\mathbb{E}_{x\sim P_{\mathrm{data}}}\left[\log p_\theta(x)\right].
\]
但在实际问题中，我们并不知道真实分布 \(P_{\mathrm{data}}\) 的具体表达式，只知道一组样本
\[
x_1,x_2,\dots,x_N.
\]
于是就用样本平均来近似这个期望：
\[
\mathbb{E}_{x\sim P_{\mathrm{data}}}[\log p_\theta(x)]
\approx
\frac{1}{N}\sum_{i=1}^N \log p_\theta(x_i).
\]
因此，优化目标就变成
\[
\max_\theta \frac{1}{N}\sum_{i=1}^N \log p_\theta(x_i).
\]


因为对数函数是单调递增的，最大化
\[
\frac{1}{N}\sum_{i=1}^N \log p_\theta(x_i)
\]
等价于最大化
\[
\prod_{i=1}^N p_\theta(x_i).
\]
后者正是样本的似然函数：
\[
L(\theta)=\prod_{i=1}^N p_\theta(x_i).
\]
因此，最优参数可以写成
\[
\hat{\theta}_{\mathrm{MLE}}
=
\arg\max_\theta \prod_{i=1}^N p_\theta(x_i),
\]
或者等价地写成
\[
\hat{\theta}_{\mathrm{MLE}}
=
\arg\max_\theta \sum_{i=1}^N \log p_\theta(x_i).
\]

这说明最小化
\[
D_{\mathrm{KL}}(P_{\mathrm{data}}\|P_\theta)
\]
的结果，就是最大似然估计（MLE）。
